% ============================================================
% preambulo-revista-contrato.tex
% ------------------------------------------------------------
% CONTRATO DE REVISTA VERSIÓN 4
% ------------------------------------------------------------
% v4 (12/11/2026): B5, EL PDF DEL NÚMERO (5): TAPA, PÁGINA LEGAL,
% SUMARIO Y LOS ARTÍCULOS INCLUIDOS CON pdfpages. LO EMITE
% numero-maestro-latex.xsl EN latex/maestro-numero.tex.
% v3 (11/11/2026): A1.1b-3, LOS IDIOMAS POR CÓDIGO (4.0):
% \gbselectlanguage{es} Y gbresumen{en}{…} TOMAN EL CÓDIGO Y BUSCAN
% EL NOMBRE DE babel EN \gbbabel@<código>, QUE DEFINE
% preambulo-numero.tex DESDE EL CATÁLOGO DE IDIOMAS.
% v2 (07/11/2026): A1.1b, LA PRIMERA PÁGINA SOBRE EL FRONT NUEVO
% (4.2 A 4.6, CON LOS ÍCONOS DE COMPARTIR: \gbcompartir, \gbicono),
% \articuloautorescorto Y \articuloeid; SALEN textpos Y
% \gbbloquex. LO EMITE jats-a-latex-articulo.xsl.
% v1 (07/11/2026): ETAPA A1 DE LAS SALIDAS DE REVISTAS
% (claude/salidas-revistas.md). REEMPLAZA LO QUE
% m_XML.ObtenerPreambuloEmbebido DEFINÍA DENTRO DE GAMBAS PARA EL
% CAMINO VIEJO, QUE SIGUE ALLÁ HASTA EL 2.0.0.
% ------------------------------------------------------------
% DEFINE TODO LO QUE jats-a-latex-articulo.xsl (QUE IMPORTA
% jats-to-latex.xsl) PUEDE EMITIR EN EL CUERPO
% DE UN ARTÍCULO. LO COMÚN CON LIBROS VIENE DE
% preambulo-contrato-comun.tex; ACÁ VA LO DE REVISTA: LA
% BIBLIOGRAFÍA DE UN ARTÍCULO Y LAS MACROS DE SU PRIMERA PÁGINA.
% COMO EN LIBROS, ES LA CAPA QUE SE CONSERVA CUANDO SE REEMPLAZA LA
% ESTÉTICA: SI FALTA UNA DEFINICIÓN, EL ARTÍCULO NO COMPILA; SI ESTÁ
% DISTINTA, COMPILA Y SALE MAL.
% ------------------------------------------------------------
% ORDEN DE LAS CAPAS (main-{nombre}.tex, LO ESCRIBE
% engine/generar_pdf_revista.sh):
%   preambulo-numero.tex            LOS VALORES DEL NÚMERO, DE LA BASE
%   preambulo-revista-contrato.tex  ESTE ARCHIVO
%   preambulo-revista-estetica.tex  EL DISEÑO DE LA REVISTA
% ------------------------------------------------------------
% DEPENDE DE preambulo-numero.tex, QUE VA ANTES:
%   \gbidiomas \gbidiomaprincipal \ifgbcodigocolor   (CONTRATO COMÚN)
%   \gbbabel@<código>                                EL NOMBRE DE babel
%                                                    DE CADA CÓDIGO (4.0)
%   \estandar                                        EL ESTILO DE CITA
%   \gbarchivobib                                    EL .bib
% ============================================================

\def\gbContratoRevistaVersion{4}

% ============================================================
% \gbRequiereContratoRevista{N}
% PROPÓSITO : EL main-{nombre}.tex DECLARA QUÉ VERSIÓN NECESITA.
%             UN PREÁMBULO VIEJO FRENA EN LA PRIMERA LÍNEA CON UN
%             MENSAJE ÚTIL, COMO EN LIBROS (\gbRequiereContrato)
% PARÁMETRO : N — versión requerida
% ============================================================
\newcommand{\gbRequiereContratoRevista}[1]{%
  \ifnum#1>\gbContratoRevistaVersion\relax
    \PackageError{gbpublisher}{%
      Este preambulo implementa el contrato de revista
      \gbContratoRevistaVersion\space y el articulo requiere el #1}{%
      Actualizar preambulo-revista-contrato.tex.}%
  \fi}

% ============================================================
% LO COMÚN CON LIBROS
% ------------------------------------------------------------
% CORTE Y NOTA EN TÍTULOS, TACHADO Y SUBRAYADO, PAQUETES, babel,
% csquotes, CIFRAS, CUADROS, BLOQUES, LLAMADAS DE ATENCIÓN Y
% CÓDIGO. VA PRIMERO: biblatex PIDE babel Y csquotes ANTES.
% ============================================================
% ============================================================
% preambulo-contrato-comun.tex
% ------------------------------------------------------------
% CONTRATO COMÚN VERSIÓN 1
% ------------------------------------------------------------
% v1 (07/11/2026): SALE DE preambulo-contrato.tex, QUE ERA DE
% LIBROS. ES LO QUE EMITEN LAS DOS HOJAS DE LaTeX,
% docbook-to-latex.xsl (LIBROS) Y jats-to-latex.xsl (REVISTAS), Y
% NO DEPENDE DE LA CLASE NI DEL CAPÍTULO: CORTE Y NOTA EN
% TÍTULOS, TACHADO Y SUBRAYADO, LOS PAQUETES DEL CONTRATO, babel,
% csquotes, CIFRAS, CUADROS, BLOQUES (CITA, EPÍGRAFE, RECUADROS,
% VERSO, CONVERSACIÓN), LLAMADAS DE ATENCIÓN Y CÓDIGO. EL TEXTO
% DE CADA APARTADO SE MUDÓ SIN CAMBIOS, SALVO EL CONTADOR
% gbcodigo (VER «8. CÓDIGO»).
% LO CARGA, PRIMERO, EL CONTRATO DE CADA TIPO DE PROYECTO:
%   libro    preambulo-contrato.tex
%   revista  preambulo-revista-contrato.tex
% UN CAMBIO ACÁ LLEGA A LOS DOS: SE PRUEBA CON UN LIBRO Y CON
% UNA REVISTA ANTES DE SUBIR LA VERSIÓN.
% ------------------------------------------------------------
% DEPENDE DE LAS DEFINICIONES DEL PROYECTO, QUE VAN ANTES:
%   \gbidiomas          LOS IDIOMAS DE babel
%   \ifgbcodigocolor    EL CÓDIGO EN COLOR O EN GRISES (SI NO ESTÁ
%                       DEFINIDO, EN COLOR)
% ============================================================

\def\gbContratoComunVersion{1}

% ============================================================
% CORTE DE LÍNEA EN TÍTULOS (CONTRATO 6, SC-28)
% ------------------------------------------------------------
% El editor marca en la base dónde cortar un título o un
% subtítulo. El generador lo traduce a \gbCorteTitulo SOLO en
% la versión compuesta del título: la que va a la apertura de
% la pieza o de la parte y a la portada. El sumario, el folio y
% los marcadores reciben siempre la versión plana, sin la
% macro: por eso las macros de título de este contrato llevan
% las dos.
%
% \gbtitulo Y \gbsubtitulo VAN COMPUESTOS Y TAMBIÉN LLEGAN A
% pdftitle Y pdfsubject. AHÍ hyperref PASA EL VALOR A CADENA
% PDF, Y UN \\ NO TIENE LUGAR: LA DESACTIVACIÓN LO CAMBIA POR
% UN ESPACIO. VALE ESPACIO Y NO NADA PORQUE EL GENERADOR ESCRIBE
% «Uno\gbCorteTitulo dos», Y EL ESPACIO QUE SIGUE A LA MACRO LO
% CONSUME TeX AL LEER EL NOMBRE.
%
% hyperref SE CARGA EN LA ÚLTIMA CAPA, Y \pdfstringdefDisableCommands
% NO EXISTE ANTES: SE REGISTRA EN EL GANCHO DE CARGA DEL PAQUETE,
% QUE CORRE APENAS hyperref TERMINA DE CARGARSE Y ANTES DEL
% \hypersetup{pdftitle=...} DE ESA MISMA CAPA. EN \AtBeginDocument
% SERÍA TARDE: pdftitle YA ESTARÍA CONVERTIDO.
% ============================================================
\newcommand{\gbCorteTitulo}{\\}
\AddToHook{package/hyperref/after}{%
  \pdfstringdefDisableCommands{\def\gbCorteTitulo{ }}}

% ============================================================
% NOTA AL PIE EN UN TÍTULO DE SECCIÓN (CONTRATO 7, SC-29)
% ------------------------------------------------------------
% \gbNotaTitulo{TEXTO}: LA EMITE docbook-to-latex.xsl EN EL
% ARGUMENTO OBLIGATORIO DE \section Y SUS NIVELES, NUNCA EN EL
% OPCIONAL: EL SUMARIO, EL FOLIO Y LOS MARCADORES VAN SIN NOTA.
%
% CON LA OPCIÓN calcwidth, titlesec COMPONE EL TÍTULO DOS VECES,
% UNA PARA MEDIRLO, Y UN \footnote DIRECTO AVANZA DOS VECES EL
% CONTADOR: LAS NOTAS SALEN 2, 3, 5 (VERIFICADO, GV-70). EL
% MANUAL DE titlesec (§2.9) DA \iftitlemeasuring{MIDE}{COMPONE}
% PARA ESTE CASO: AL MEDIR NO SE EMITE NADA.
%
% SI LA ESTÉTICA NO CARGA titlesec, \iftitlemeasuring NO EXISTE:
% SE PROVEE AL EMPEZAR EL DOCUMENTO, CUANDO LA ESTÉTICA YA SE
% CARGÓ, Y SIEMPRE COMPONE. \providecommand NO TOCA LA DE
% titlesec.
%
% LOS PÁRRAFOS DE LA NOTA LLEGAN SEPARADOS CON \endgraf: EL
% ARGUMENTO DE UN TÍTULO NO ADMITE \par.
% ============================================================
\newcommand{\gbNotaTitulo}[1]{%
  \iftitlemeasuring{}{\footnote{#1}}}
\AtBeginDocument{%
  \providecommand{\iftitlemeasuring}[2]{#2}}

% ============================================================
% TACHADO Y SUBRAYADO (CONTRATO 8, SC-30)
% ------------------------------------------------------------
% \gbTachado{TEXTO} Y \gbSubrayado{TEXTO}: LOS EMITE
% docbook-to-latex.xsl PARA ~~x~~ Y [x]{.underline}.
%
% ulem Y NO \underline: \underline ARMA UNA CAJA Y NO CORTA LA
% LÍNEA, Y LaTeX NO TIENE TACHADO PROPIO. [normalem] ES
% OBLIGATORIO: SIN ÉL, ulem CAMBIA \emph POR SUBRAYADO EN TODO
% EL LIBRO.
%
% \DeclareRobustCommand PORQUE PUEDEN QUEDAR EN UN TÍTULO DE
% SECCIÓN, QUE VIAJA AL SUMARIO, Y LOS COMANDOS DE ulem SON
% FRÁGILES. EN LOS MARCADORES DEL PDF SE QUEDA EL TEXTO SOLO.
% ============================================================
\usepackage[normalem]{ulem}
\DeclareRobustCommand{\gbTachado}[1]{\sout{#1}}
\DeclareRobustCommand{\gbSubrayado}[1]{\uline{#1}}
\AddToHook{package/hyperref/after}{%
  \pdfstringdefDisableCommands{%
    \def\gbTachado#1{#1}%
    \def\gbSubrayado#1{#1}}}

% ============================================================
% 1. PAQUETES QUE EL CONTRATO NECESITA
% ------------------------------------------------------------
% SOLO LO QUE HACE FALTA PARA QUE UN FRAGMENTO COMPILE. LA
% ESTÉTICA CARGA LO SUYO APARTE. TODOS ESTÁN EN TeX Live Y
% NINGUNO PIDE --shell-escape.
% ============================================================
\usepackage{etoolbox}
\usepackage{graphicx}
\usepackage{booktabs}
\usepackage{array}
% fancyvrb: EL Verbatim CON commandchars DEL CÓDIGO COLOREADO (SC-42)
\usepackage{fancyvrb}
\usepackage{ragged2e}
\usepackage{xcolor}
\usepackage[most]{tcolorbox}
% SEPARADOR ORNAMENTAL \froufrou (SC-36). CARGA tikz Y fourier-orns
\usepackage{froufrou}
% VERSO (SC-49): REDEFINE verse Y TRAE \vin, \\! Y \stanzaskip. EN
% TeX Live 2023 (MINT) ES LA 2.4a; LO QUE SE USA EXISTE DESDE LA 2.2
\usepackage{verse}

% ------------------------------------------------------------
% IDIOMA
% ------------------------------------------------------------
% babel VA EN EL CONTRATO Y NO EN LA ESTÉTICA, POR DOS RAZONES.
% LA PRIMERA ES DE ORDEN: biblatex PIDE CARGARSE DESPUÉS DE
% babel Y csquotes, Y LA ESTÉTICA VA DESPUÉS DE ESTA CAPA.
% LA SEGUNDA ES DE FONDO: EL IDIOMA DECIDE LAS COMILLAS QUE
% PRODUCE \enquote, LA PARTICIÓN DE PALABRAS Y LAS CADENAS DE
% LA BIBLIOGRAFÍA. SON GARANTÍAS DEL CONTRATO, NO DECISIONES
% DE DISEÑO: UNA EDITORIAL QUE REEMPLAZA LA ESTÉTICA NO DEBE
% PODER CAMBIAR EL IDIOMA DEL LIBRO.
\usepackage[\gbidiomas]{babel}
\frenchspacing

% COMILLAS SEGÚN EL IDIOMA ACTIVO: EN ESPAÑOL DA « ».
% EL XSLT EMITE \enquote Y NUNCA COMILLAS LITERALES.
\usepackage[autostyle=true]{csquotes}

% ------------------------------------------------------------
% CIFRAS Y CANTIDADES (CONTRATO 15, SC-64, SC-65)
% ------------------------------------------------------------
% LA NORMA INTERNACIONAL (SI Brochure §5.4): PUNTO DECIMAL,
% GRUPOS DE A TRES CON UN ESPACIO DESDE CINCO DÍGITOS, LA UNIDAD
% Y EL PORCENTAJE SEPARADOS DE LA CIFRA, EL GRADO PEGADO.
% detect-all NO SE QUITA: SIN ELLA siunitx COMPONE LA CIFRA EN
% MODO MATEMÁTICO, DONDE babel-spanish PASA EL PUNTO A COMA, Y
% LOS DECIMALES SALEN «12,5» SIN AVISO (GV-123). EN siunitx 3 ES
% UNA OPCIÓN DEPRECADA: DEJA UN AVISO INFORMATIVO EN EL LOG.
% per-mode=symbol: «48 km/h» EN UNA LÍNEA; LA FRACCIÓN VERTICAL
% AGRANDABA LA LÍNEA (DECISIÓN DE ALBERTO).
% EL FORMATEADOR DE cifras.lua APLICA LAS MISMAS REGLAS PARA EL
% HTML Y EL EPUB: SI CAMBIA ESTE \sisetup, CAMBIA EL FILTRO.
\usepackage{siunitx}
\sisetup{output-decimal-marker={.},
  group-separator={\hspace{0.15em}},
  group-minimum-digits=5,
  table-text-alignment=center,
  detect-all,
  per-mode=symbol}
% siunitx NO DEFINE EL EURO (GV-123). LA FUENTE TIENE QUE TENER €
\DeclareSIUnit\euro{€}

% ------------------------------------------------------------
% CUADROS (CONTRATO 16, SC-66)
% ------------------------------------------------------------
% EL ESTILO DE LA CASA: LETRA SIN SERIFA, FILETE ENTRE TODAS LAS
% FILAS Y ENTRE TODAS LAS COLUMNAS, ENCABEZADO EN NEGRITA, PIE
% DEBAJO. LOS FILETES Y LA NEGRITA LOS ESCRIBE EL XSLT; ACÁ, LO
% QUE UNA EDITORIAL PUEDE CAMBIAR SIN TOCARLO. EL CUERPO Y LA
% SEPARACIÓN DE COLUMNAS LOS ELIGE EL MAQUETADOR POR CUADRO
% (cuerpo: normal|chico, separacion: normal|estrecha). LA LETRA
% Y EL CUERPO VALEN PARA TODO EL FLOTANTE, EL PIE INCLUIDO.
% GEMELAS: m_XML.ObtenerPreambuloEmbebido (REVISTAS).
\newcommand{\gbCuadroLetra}{\sffamily}
\newcommand{\gbCuadroCuerpoNormal}{\footnotesize}
\newcommand{\gbCuadroCuerpoChico}{\scriptsize}
\newcommand{\gbCuadroSepNormal}{4pt}
\newcommand{\gbCuadroSepEstrecha}{2pt}
% UN CUADRO CORTADO (cortes:) VA EN UN FLOTANTE POR TRAMO: EL PRIMERO
% TERMINA CON \gbCuadroContinua, A LA DERECHA; LOS QUE SIGUEN EMPIEZAN
% CON \gbCuadroViene, A LA IZQUIERDA, Y REPITEN EL ENCABEZADO. EL PIE,
% SOLO EN EL ÚLTIMO. EL XSLT LOS PONE EN BASTARDILLA.
\newcommand{\gbCuadroContinua}{Continúa en la página siguiente}
\newcommand{\gbCuadroViene}{Viene de la página anterior}
% EL CUADRO APAISADO (rotar: sí): sidewaystable, UNA PÁGINA POR TRAMO
\usepackage{rotating}

% ============================================================
% 6. BLOQUES
% ============================================================

% CITA EN BLOQUE CON ATRIBUCIÓN OPCIONAL.
% NO ES UN EPÍGRAFE: DOCBOOK LOS DISTINGUE Y EL XSLT TAMBIÉN.
\newenvironment{gbcita}
  {\begin{quote}\small}
  {\end{quote}}
\newcommand{\gbatribucion}[1]{\par\vspace{2pt}\hfill\textup{\small#1}\par}

% EPÍGRAFE (SC-35). BLOQUE A LA DERECHA, 60 % DE LA COLUMNA, \small,
% SIN CORTE DE PALABRA. #1 A LA IZQUIERDA; FILETE DE 0,6 PT; #2 A LA
% DERECHA. CON #2 VACÍO NO HAY FILETE. ENTRE PÁRRAFOS DE #1, MEDIA LÍNEA. #1 PUEDE TENER VARIOS PÁRRAFOS
% (\newcommand SIN ASTERISCO ADMITE \par). #2 PUEDE SER UNA CITA: EL
% XSLT RESUELVE UN <biblioref> A \textcite.
% GEMELA DE LA DEFINICIÓN DE m_XML.ObtenerPreambuloEmbebido (REVISTAS):
% UN CAMBIO EN UNA VA EN LAS DOS.
\newcommand{\gbepigrafe}[2]{%
  \par\vspace{.5\baselineskip}%
  {\noindent\hfill
   \begin{minipage}{.6\linewidth}%
     \small\parindent=0pt\parskip=.5em
     \hyphenpenalty=10000\exhyphenpenalty=10000
     \raggedright #1\par
     \if\relax\detokenize{#2}\relax\else
       \parskip=0pt\vspace{1ex}\hrule height 0.6pt\vspace{1ex}%
       {\raggedleft #2\par}%
     \fi
   \end{minipage}\par}%
  \vspace{.5\baselineskip}}

% RECUADRO LATERAL. UN ARGUMENTO: info | supplementary.
\newenvironment{gbsidebar}[1]
  {\begin{tcolorbox}[colback=black!4,colframe=black!4,
     boxrule=0pt,arc=2pt,left=6pt,right=6pt,top=4pt,bottom=4pt,
     before skip=6pt,after skip=6pt]\small\sffamily}
  {\end{tcolorbox}}

% ============================================================
% RECUADROS (CONTRATO 10, SC-41)
% ------------------------------------------------------------
% gbRecuadro: FONDO GRIS AL 7 %, FILETE DE 1 PT Y ESQUINAS DE
% 5 PT, LOS DOS EN gbcolorbase.
% gbRecuadroBarra{BARRA}: FONDO GRIS AL 7 %, SIN FILETE Y CON
% ESQUINAS RECTAS; ARRIBA, UNA BARRA DE UNA PICA EN gbcolorbase
% CON EL TEXTO EN SANS NEGRITA Y MAYÚSCULAS, BLANCO Y CENTRADO
% EN ALTURA: EL MISMO RÓTULO QUE EL TIPO DE ARTÍCULO DE LA
% CABECERA. \MakeUppercase Y NO upper-case() EN EL XSLT: CONSERVA
% LA BASTARDILLA Y DEJA INTACTA LA CLAVE DE UN \autocite (MEDIDO).
% EL TEXTO ES EL DEL CUERPO: breakable PARA QUE UN RECUADRO
% LARGO CRUCE DE PÁGINA; parbox=false PARA QUE LOS PÁRRAFOS
% TENGAN LA SANGRÍA DEL CUERPO (EN MODO MINIPÁGINA QUEDARÍAN
% PEGADOS Y SIN SANGRÍA); EL PRIMERO VA SIN SANGRÍA.
% SIN NOTAS AL PIE: EN UN tcolorbox QUEDAN ADENTRO, CON LETRA
% (MEDIDO); EL FILTRO recuadros.lua LAS FRENA.
% EL COLOR LO DEFINE LA ESTÉTICA (\definecolor{gbcolorbase});
% ESTE \providecolor ES EL RESPALDO DE UNA ESTÉTICA QUE NO LO
% DECLARE. GEMELOS EN m_XML.ObtenerPreambuloEmbebido (REVISTAS):
% UN CAMBIO EN UNO VA EN LOS DOS.
% ============================================================
\providecolor{gbcolorbase}{HTML}{2D5A8E}
\tcbset{
  gbrecuadro/.style={breakable, parbox=false,
    colback=black!7, colframe=gbcolorbase, boxrule=1pt, arc=5pt,
    boxsep=0pt, left=8pt, right=8pt, top=8pt, bottom=8pt,
    before skip=\baselineskip, after skip=\baselineskip,
    before upper=\noindent},
  gbrecuadrobarra/.style={breakable, parbox=false, sharp corners,
    colback=black!7, colframe=black!7, boxrule=0pt,
    colbacktitle=gbcolorbase, coltitle=white, titlerule=0pt,
    boxsep=0pt, left=8pt, right=8pt, top=8pt, bottom=8pt,
    toptitle=0pt, bottomtitle=0pt, lefttitle=8pt, righttitle=8pt,
    before skip=\baselineskip, after skip=\baselineskip,
    before upper=\noindent}}
\newenvironment{gbRecuadro}
  {\begin{tcolorbox}[gbrecuadro]}
  {\end{tcolorbox}}
\newenvironment{gbRecuadroBarra}[1]
  {\begin{tcolorbox}[gbrecuadrobarra,
     title={\parbox[c][1pc][c]{\linewidth}{\sffamily\bfseries\MakeUppercase{#1}}}]}
  {\end{tcolorbox}}

% ============================================================
% VERSO (CONTRATO 14, SC-49)
% ------------------------------------------------------------
% EL ENTORNO verse DEL PAQUETE verse. EL XSLT TERMINA CADA VERSO
% EN \\ Y EL ÚLTIMO DE CADA ESTROFA EN \\!, QUE DEJA \stanzaskip
% (0,75 DE LÍNEA) ENTRE ESTROFAS; EL ÚLTIMO DEL POEMA VA SIN NADA.
% LA SANGRÍA ES \vin (\vgap, 1,5 em) UNA VEZ POR NIVEL; UN VERSO
% QUE NO ENTRA SIGUE A \vindent (3 em). CADA VERSO EMPIEZA CON
% \relax, PARA QUE UN [ O UN * DEL TEXTO NO SE LEA COMO ARGUMENTO
% DE \\ (GV-93).
% gbverso: EL POEMA, EN BASTARDILLA, A \leftmargini DEL MARGEN.
% gbversoepigrafe: EL VERSO DE UN EPÍGRAFE, DENTRO DEL PRIMER
% ARGUMENTO DE \gbepigrafe: LA LETRA DEL EPÍGRAFE, SIN MARGEN
% PROPIO Y SIN EL \topsep DE LA LISTA, QUE SEPARARÍA EL ÚLTIMO
% VERSO DEL FILETE. LOS DOS SALEN DE \@listi, QUE \list LEE AL
% ABRIR: SE REDEFINE LOCAL AL ENTORNO. MEDIDO: DEL CUERPO AL
% PRIMER VERSO Y DEL ÚLTIMO VERSO A LA ATRIBUCIÓN, LA MISMA
% DISTANCIA QUE EN UN EPÍGRAFE EN PROSA (GV-94). \parsep E
% \itemsep LOS PONE verse DESPUÉS (\stanzaskip ENTRE ESTROFAS).
% GEMELOS EN m_XML.ObtenerPreambuloEmbebido (REVISTAS): UN CAMBIO
% EN UNOS VA EN LOS OTROS.
% ============================================================
\newenvironment{gbverso}
  {\begin{verse}\itshape}
  {\end{verse}}
\makeatletter
\newenvironment{gbversoepigrafe}
  {\def\@listi{\leftmargin\z@\topsep\z@\partopsep\z@
     \itemsep\z@\parsep\z@}%
   \begin{verse}}
  {\end{verse}}
\makeatother

% ============================================================
% CONVERSACIÓN (CONTRATO 12, SC-43)
% ------------------------------------------------------------
% gbconversacion: EL qandaset role="conversacion" ENTERO. SANS,
% \small; TODO EL BLOQUE CORRIDO 14 PT A LA IZQUIERDA (\leftskip), EL
% MISMO MARGEN QUE LA CITA (quote, EN preambulo-estetica.tex), SIN
% MARGEN A LA DERECHA; SIN SANGRÍA DE PRIMERA LÍNEA. \leftskip Y NO
% UNA LISTA COMO LA CITA: DENTRO DE UNA LISTA EL ESPACIO ENTRE PÁRRAFOS
% ES \parsep, NO \parskip. LAS NOTAS NO LO HEREDAN: \@parboxrestore
% LO PONE EN CERO. UN ESPACIO CHICO (\smallskipamount) ENTRE
% TURNOS Y ENTRE LOS PÁRRAFOS DE UN TURNO, QUE SIN SANGRÍA ES LO
% ÚNICO QUE LOS SEPARA; \medskip AL ABRIR Y AL CERRAR, QUE DA LA
% SENSACIÓN DE BLOQUE.
% \gbetiqueta{ETIQUETA}: ABRE EL PRIMER PÁRRAFO DEL TURNO, EN
% LÍNEA. LA HOJA LA ESCRIBE YA EN MAYÚSCULAS; PESO NORMAL.
% LA PREGUNTA NO QUEDA SOLA AL PIE: LA HOJA ESCRIBE \nopagebreak[4]
% DESPUÉS DE CADA UNA.
% GEMELO DE LA DEFINICIÓN DE m_XML.ObtenerPreambuloEmbebido
% (REVISTAS): UN CAMBIO EN UNA VA EN LAS DOS.
% ============================================================
\newenvironment{gbconversacion}
  {\par\medskip
   \begingroup\sffamily\small
   \setlength{\leftskip}{14pt}%
   \setlength{\parindent}{0pt}%
   \setlength{\parskip}{\smallskipamount}}
  {\par\endgroup\medskip}
\newcommand{\gbetiqueta}[1]{#1\enspace}

% ============================================================
% 7. LLAMADAS DE ATENCIÓN
% ============================================================
% CADA LLAMADA DE ATENCIÓN ES UN tcolorbox CON TÍTULO FIJO.
% NO SE USA UN .style DE pgfkeys PARAMETRIZADO: EL #1 DE UN
% \newcommand DENTRO DE UN .style LO CONSUME pgfkeys Y DA
% "You already have nine parameters", QUE NO DICE NADA SOBRE
% LA CAUSA. CON UN ENTORNO POR CLASE EL PROBLEMA NO EXISTE.
\tcbset{gbadmon/.style={
  colback=black!3, colframe=black!25,
  boxrule=0.4pt, arc=1pt,
  left=6pt, right=6pt, top=4pt, bottom=4pt,
  before skip=6pt, after skip=6pt,
  fonttitle=\bfseries\sffamily\small}}

\newenvironment{gbnota}
  {\begin{tcolorbox}[gbadmon,title={Nota}]\small}{\end{tcolorbox}}
\newenvironment{gbconsejo}
  {\begin{tcolorbox}[gbadmon,title={Sugerencia}]\small}{\end{tcolorbox}}
\newenvironment{gbadvertencia}
  {\begin{tcolorbox}[gbadmon,title={Advertencia}]\small}{\end{tcolorbox}}
\newenvironment{gbimportante}
  {\begin{tcolorbox}[gbadmon,title={Importante}]\small}{\end{tcolorbox}}
\newenvironment{gbprecaucion}
  {\begin{tcolorbox}[gbadmon,title={Precaución}]\small}{\end{tcolorbox}}

% ============================================================
% 8. CÓDIGO (CONTRATO 11, SC-42)
% ------------------------------------------------------------
% EL COLOR NO LO PONE LaTeX: LO PONE EL RESALTADOR DE PANDOC
% (skylighting) EN UN PASO PREVIO (colorear_codigo.sh), QUE
% ESCRIBE CADA LÍNEA CON LAS MACROS DE TOKEN (\KeywordTok{for}
% Y LAS DEMÁS). ACÁ SOLO SE DEFINE CÓMO SE VE CADA UNA. NI
% listings NI minted: listings NO CONOCE JavaScript, JSON, YAML
% NI Lua (MEDIDO), Y minted EXIGE --shell-escape Y Pygments.
%
% gbCodigo{RÓTULO}: CAJA OSCURA CON CABECERA, EL RÓTULO DEL
% LENGUAJE ARRIBA Y LAS LÍNEAS EN UN Verbatim CON commandchars.
% breakable: UN BLOQUE LARGO CRUZA DE PÁGINA.
% \gbnl{N}: NÚMERO DE LÍNEA, AL MARGEN Y EN GRIS.
% \gbnc: CONTINUACIÓN DE UN CORTE MANUAL, SIN NÚMERO.
% \gbret: EL ↩ DEL CORTE MANUAL, EN EL GRIS DE LOS NÚMEROS.
% \gbPieCodigo{ID}{TEXTO}: EL PIE DE UN LISTADO, DEBAJO DEL
% BLOQUE COMO EN LAS FIGURAS: «Código 2.3.», CON EL CONTADOR
% gbcodigo (DECLARADO CON LOS DE LA NUMERACIÓN POR PIEZA).
%
% \ifgbcodigocolor: VERDADERO, PALETA Gruvbox OSCURA (EL TEMA
% DE LA APLICACIÓN); FALSO, FONDO BLANCO, TEXTO NEGRO Y SIN
% COLOR, PARA IMPRENTA. LA BANDERA LA ESCRIBE EL GENERADOR EN
% LAS PRESTACIONES DE SALIDA (estado «publicado»); SI NO LLEGA,
% VA EL COLOR. LA NEGRITA Y LA BASTARDILLA DE LOS TOKENS SE
% CONSERVAN EN LOS DOS CASOS.
%
% LA PALETA ES LA DE .hidden/themes/gruvbox.theme. GEMELOS: EL
% CSS DEL CÓDIGO EN gbpublisher.css Y EN jats-to-html.xsl, Y
% ESTA MISMA SECCIÓN EN m_XML.ObtenerPreambuloEmbebido
% (REVISTAS). UN CAMBIO DE COLOR VA EN LOS CUATRO.
% ============================================================
\expandafter\ifx\csname ifgbcodigocolor\endcsname\relax
  \expandafter\newif\csname ifgbcodigocolor\endcsname
  \gbcodigocolortrue
\fi

\ifgbcodigocolor
  \definecolor{gbcodfondo}{HTML}{282828}
  \definecolor{gbcodmarco}{HTML}{282828}
  \definecolor{gbcodtexto}{HTML}{EBDBB2}
  \definecolor{gbcodcabeza}{HTML}{3C3836}
  \definecolor{gbcodrotulo}{HTML}{A89984}
  \definecolor{gbcodnumero}{HTML}{7C6F64}
\else
  \definecolor{gbcodfondo}{gray}{1}
  \definecolor{gbcodmarco}{gray}{0}
  \definecolor{gbcodtexto}{gray}{0}
  \definecolor{gbcodcabeza}{gray}{0.9}
  \definecolor{gbcodrotulo}{gray}{0}
  \definecolor{gbcodnumero}{gray}{0.45}
\fi

% \gbDefTok{NOMBRE}{COLOR HTML}{ESTILO}: DEFINE \NOMBRETok. EL
% ESTILO ES \textbf, \textit O \relax. SIN COLOR, SOLO EL ESTILO.
\newcommand{\gbDefTok}[3]{%
  \ifgbcodigocolor
    \expandafter\newcommand\csname #1Tok\endcsname[1]{\textcolor[HTML]{#2}{#3{##1}}}%
  \else
    \expandafter\newcommand\csname #1Tok\endcsname[1]{#3{##1}}%
  \fi}
% LAS 31 CLASES DE skylighting
\gbDefTok{Normal}{EBDBB2}{\relax}
\gbDefTok{Variable}{EBDBB2}{\relax}
\gbDefTok{Keyword}{83A598}{\textbf}
\gbDefTok{ControlFlow}{83A598}{\textbf}
\gbDefTok{DataType}{FABD2F}{\relax}
\gbDefTok{Attribute}{FABD2F}{\relax}
\gbDefTok{DecVal}{D3869B}{\relax}
\gbDefTok{BaseN}{D3869B}{\relax}
\gbDefTok{Float}{D3869B}{\relax}
\gbDefTok{Constant}{D3869B}{\relax}
\gbDefTok{Char}{B8BB26}{\relax}
\gbDefTok{String}{B8BB26}{\relax}
\gbDefTok{VerbatimString}{B8BB26}{\relax}
\gbDefTok{SpecialString}{B8BB26}{\relax}
\gbDefTok{SpecialChar}{FE8019}{\relax}
\gbDefTok{Import}{FE8019}{\relax}
\gbDefTok{Preprocessor}{FE8019}{\relax}
\gbDefTok{Comment}{928374}{\textit}
\gbDefTok{Documentation}{928374}{\textit}
\gbDefTok{Annotation}{928374}{\textit}
\gbDefTok{CommentVar}{928374}{\textit}
\gbDefTok{Information}{928374}{\textit}
\gbDefTok{RegionMarker}{928374}{\relax}
\gbDefTok{Other}{8EC07C}{\relax}
\gbDefTok{Function}{8EC07C}{\relax}
\gbDefTok{Operator}{8EC07C}{\relax}
\gbDefTok{BuiltIn}{8EC07C}{\relax}
\gbDefTok{Extension}{8EC07C}{\relax}
\gbDefTok{Warning}{FE8019}{\textbf}
\gbDefTok{Alert}{FB4934}{\textbf}
\gbDefTok{Error}{FB4934}{\relax}

\newcommand{\gbnl}[1]{\llap{\textcolor{gbcodnumero}{#1}\hspace{1.2em}}}
\newcommand{\gbnc}{\relax}
\newcommand{\gbret}{\textcolor{gbcodnumero}{\symbol{"21A9}}}

\tcbset{gbcodigo/.style={breakable, colback=gbcodfondo,
  colframe=gbcodmarco, coltext=gbcodtexto, boxrule=0.4pt, arc=3pt,
  colbacktitle=gbcodcabeza, coltitle=gbcodrotulo, titlerule=0pt,
  fonttitle=\ttfamily\scriptsize, toptitle=1pt, bottomtitle=1pt,
  boxsep=0pt, left=3em, right=6pt, top=5pt, bottom=5pt,
  lefttitle=6pt, before skip=\baselineskip, after skip=\baselineskip}}
% \VerbatimEnvironment: fancyvrb BUSCA \end{gbCodigo} Y NO
% \end{Verbatim}. ES LA FORMA DOCUMENTADA DE ENVOLVER UN Verbatim.
\newenvironment{gbCodigo}[1]
  {\VerbatimEnvironment
   \begin{tcolorbox}[gbcodigo,title={#1}]%
   \begin{Verbatim}[commandchars=\\\{\},fontsize=\footnotesize]}
  {\end{Verbatim}\end{tcolorbox}}

% EL CONTADOR DE LOS LISTADOS (SC-42). ANTES SE DECLARABA EN EL
% CONTRATO DE LIBROS CON \newcounter{gbcodigo}[chapter]; LA CLASE
% article NO TIENE chapter. EL REINICIO POR CAPÍTULO LO AGREGA EL
% CONTRATO DE LIBROS CON \counterwithin*, QUE HACE LO MISMO QUE EL
% ARGUMENTO OPCIONAL (\@addtoreset) Y NO TOCA \thegbcodigo
\newcounter{gbcodigo}

% EL PIE VA PEGADO AL BLOQUE: SE DESCUENTA PARTE DEL after skip
\newcommand{\gbPieCodigo}[2]{%
  \par\nopagebreak\vspace{-.5\baselineskip}%
  \refstepcounter{gbcodigo}\label{#1}%
  {\small\noindent\textbf{Código~\thegbcodigo.}\space #2\par}%
  \vspace{.5\baselineskip}}


% ============================================================
% 1. BIBLIOGRAFÍA
% ------------------------------------------------------------
% LOS MISMOS cita-<estándar>-config.tex QUE LIBROS (SC-57): EL
% ESTILO ES ATRIBUTO DE LA REVISTA (revista.estilo_cita). DESDE EL
% CONTRATO 21 DE LIBROS ESOS ARCHIVOS NO TRAEN citereset=chapter:
% UN ARTÍCULO NO TIENE CAPÍTULOS, Y SU PRIMERA CITA COMPLETA ES
% LA PRIMERA DEL ARTÍCULO.
% EL ÍNDICE DE AUTORES ES DE LIBROS (esindex, preambulo-contrato.tex):
% ACÁ biblatex NO INDEXA NOMBRES.
% EL .bib LO DECLARA ESTA CAPA; \printbibliography LO ESCRIBE EL
% main-{nombre}.tex AL FINAL DEL CUERPO.
% ============================================================
\input{cita-\estandar-config.tex}
\ExecuteBibliographyOptions{indexing=false}
% ============================================================
% cita-cadenas-es.tex — CADENAS EN CASTELLANO COMUNES A TODOS LOS
% ESTILOS DE biblatex
% ------------------------------------------------------------
% SOLO LAS CADENAS QUE NI biblatex (spanish.lbx) NI EL ESTILO TRAEN:
% SON LA TRADUCCIÓN DE LOS CÓDIGOS DEL CATÁLOGO DE LA APLICACIÓN
% (TABLA cmb_biblatex: editor_type, thesis_type, pubstate) Y DE LOS
% AGREGADOS PROPIOS DE LA EDITORIAL. UNA CADENA ASÍ SOLO APARECE SI
% UNA ENTRADA LA USA: EN UN ESTILO QUE NO LA NECESITA NO CAMBIA NADA.
%
% LAS QUE spanish.lbx YA TRAE Y UN ESTILO REDEFINE (meses, url,
% phdthesis, compiler…) NO VAN ACÁ: SON DECISIONES DE CADA ESTILO Y
% QUEDAN EN SU cita-<estilo>-config.tex.
%
% LO CARGA QUIEN CARGA EL ESTILO, INMEDIATAMENTE DESPUÉS DE
% cita-<estilo>-config.tex (\NewBibliographyString EXIGE biblatex YA
% CARGADO):
%   libros  — preambulo-contrato.tex (PDF) Y EL DOCUMENTO MÍNIMO DE
%             preparar_citas_libro.sh (PASADA DE tex4ht), POR NOMBRE
%   revistas — m_XML, CON RUTA ABSOLUTA A ~/.gbpublisher/latex/
% NO LO CARGA EL CONFIG: \input RESUELVE CONTRA LA CARPETA DE LA
% COMPILACIÓN, NO CONTRA LA DEL ARCHIVO QUE LO PIDE, Y EN REVISTAS EL
% CONFIG SE CARGA POR RUTA ABSOLUTA DESDE OTRA CARPETA.
%
% \NewBibliographyString CON UNA CADENA QUE YA EXISTE NO HACE NADA
% (biblatex.sty); CON UNA NUEVA, LA DECLARA PARA TODOS LOS IDIOMAS.
% ============================================================

% SIN biblatex NO HAY CADENAS QUE DEFINIR: ES EL CASO DE vancouver,
% QUE VA POR bibtex (cita-vancouver-config.tex). EL ARCHIVO TERMINA
% ACÁ EN LUGAR DE CORTAR LA COMPILACIÓN CON \NewBibliographyString
% INDEFINIDO
\makeatletter
\@ifpackageloaded{biblatex}{\makeatother}{\makeatother\endinput}

% ------------------------------------------------------------
% 1. FUNCIONES DE EDITOR (editortype) QUE spanish.lbx NO TRAE
% ------------------------------------------------------------
% EL CATÁLOGO OFRECE coordinator Y direction. biblatex BUSCA LA
% CADENA DEL MISMO NOMBRE, SU PLURAL (…s) Y LA FORMA by…; SIN
% ELLAS IMPRIME EL CÓDIGO TAL CUAL («coordinator», MEDIDO EN IEEE).
% organizator ES UN ALIAS HEREDADO DE ISO 690: EL CATÁLOGO USA
% organizer, QUE spanish.lbx YA TRAE.
\NewBibliographyString{coordinator}
\NewBibliographyString{coordinators}
\NewBibliographyString{bycoordinator}
\NewBibliographyString{direction}
\NewBibliographyString{directions}
\NewBibliographyString{bydirection}
\NewBibliographyString{organizator}
\NewBibliographyString{organizators}
\NewBibliographyString{byorganizator}

% ------------------------------------------------------------
% 2. TIPOS DE TESIS Y TRABAJOS FINALES (thesis_type)
% ------------------------------------------------------------
% LOS CÓDIGOS DEL CATÁLOGO, SALVO mathesis, phdthesis Y candthesis,
% QUE spanish.lbx TRAE Y CADA ESTILO REDEFINE EN SU CONFIG. LOS
% ÚLTIMOS SEIS SON ALIAS HEREDADOS DE ISO 690, QUE EL CATÁLOGO NO
% OFRECE: QUEDAN PARA UN .bib ESCRITO A MANO.
\NewBibliographyString{bachelorthesis}
\NewBibliographyString{engineerthesis}
\NewBibliographyString{specializationthesis}
\NewBibliographyString{fieldwork}
\NewBibliographyString{finalessay}
\NewBibliographyString{finalmonograph}
\NewBibliographyString{finalproject}
\NewBibliographyString{gradothesis}
\NewBibliographyString{integrativefinalwork}
\NewBibliographyString{magisterthesis}
\NewBibliographyString{masterfinalwork}
\NewBibliographyString{postgradthesis}
\NewBibliographyString{teacherthesis}
\NewBibliographyString{researchproject}
\NewBibliographyString{residencyreport}
\NewBibliographyString{shortthesis}
\NewBibliographyString{technicalreport}
\NewBibliographyString{undergraduatethesis}
\NewBibliographyString{ingenieriathesis}
\NewBibliographyString{licenciaturathesis}
\NewBibliographyString{maestriathesis}
\NewBibliographyString{posgradothesis}
\NewBibliographyString{profesoradothesis}
\NewBibliographyString{tifthesis}

% ------------------------------------------------------------
% 3. ESTADOS DE PUBLICACIÓN (pubstate) Y FORMAS DE EDICIÓN
% ------------------------------------------------------------
\NewBibliographyString{mimeograph}
\NewBibliographyString{photocopy}
\NewBibliographyString{digitalprinting}
\NewBibliographyString{edauthor}
% self-published ES EL CÓDIGO DEL CATÁLOGO PARA LA EDICIÓN DE AUTOR (SIN
% CADENA, biblatex LO IMPRIMÍA TAL CUAL: MEDIDO EN VERONA E IEEE). edauthor
% QUEDA PARA UN .bib ESCRITO A MANO. EL GUION EN LA CLAVE ES VÁLIDO (MEDIDO)
\NewBibliographyString{self-published}
\NewBibliographyString{documentjob}
\NewBibliographyString{origpubbare}

% ------------------------------------------------------------
% 4. CINE: DIRECCIÓN, GUION Y ELENCO
% ------------------------------------------------------------
\NewBibliographyString{dirigida}
\NewBibliographyString{dirigidas}
\NewBibliographyString{bydirigida}
\NewBibliographyString{escrita}
\NewBibliographyString{escritas}
\NewBibliographyString{byescrita}
\NewBibliographyString{elenco}
\NewBibliographyString{elencos}
\NewBibliographyString{byelenco}

% ------------------------------------------------------------
% 5. «Y SIGUIENTE», «Y SIGUIENTES»: \sq Y \sqq EN UNA postnote
% ------------------------------------------------------------
\NewBibliographyString{sq}
\NewBibliographyString{sqq}
\providecommand*{\sq}{\bibstring{sq}}
\providecommand*{\sqq}{\bibstring{sqq}}

% ------------------------------------------------------------
% 6. EL TEXTO EN CASTELLANO
% ------------------------------------------------------------
\DefineBibliographyStrings{spanish}{%
  coordinator          = {coord\adddot},
  coordinators         = {coords\adddot},
  bycoordinator        = {coord\adddot\addspace por},
  direction            = {dir\adddot},
  directions           = {dirs\adddot},
  bydirection          = {dir\adddot\addspace por},
  organizator          = {org\adddot},
  organizators         = {orgs\adddot},
  byorganizator        = {org\adddot\addspace por},
  bachelorthesis       = {Tesis de Licenciatura},
  engineerthesis       = {Tesis de Ingeniería},
  specializationthesis = {Tesis de Especialización},
  fieldwork            = {Trabajo de Campo},
  finalessay           = {Ensayo Final},
  finalmonograph       = {Monografía Final de Carrera},
  finalproject         = {Proyecto Final},
  gradothesis          = {Tesis de Grado},
  integrativefinalwork = {Trabajo Integrador Final (TIF)},
  magisterthesis       = {Tesis de Magíster},
  masterfinalwork      = {Trabajo Final de Maestría},
  postgradthesis       = {Tesis de Posgrado},
  teacherthesis        = {Tesis de Profesorado},
  researchproject      = {Proyecto de Investigación},
  residencyreport      = {Informe Final de Residencia},
  shortthesis          = {Tesina},
  technicalreport      = {Trabajo Final Técnico},
  undergraduatethesis  = {Trabajo Final de Grado},
  ingenieriathesis     = {Tesis de Ingeniería},
  licenciaturathesis   = {Tesis de Licenciatura},
  maestriathesis       = {Tesis de Maestría},
  posgradothesis       = {Tesis de Posgrado},
  profesoradothesis    = {Tesis de Profesorado},
  tifthesis            = {Trabajo Integrador Final},
  mimeograph           = {mimeo},
  photocopy            = {fotocopia},
  digitalprinting      = {impresión digital},
  edauthor             = {edición de autor},
  self-published       = {edición de autor},
  documentjob          = {documento de trabajo},
  origpubbare          = {orig\adddotspace pub\adddotspace},
  dirigida             = {dirección},
  dirigidas            = {dirección},
  bydirigida           = {Dirección de},
  escrita              = {escrita},
  escritas             = {escrita},
  byescrita            = {escrita por},
  elenco               = {actúan:},
  elencos              = {actúan:},
  byelenco             = {actúan:},
  sq                   = {~y~sig\adddot},
  sqq                  = {~y~ss\adddot},
}

\addbibresource{\gbarchivobib}

% ============================================================
% 2. PAQUETES QUE EL CUERPO DEL ARTÍCULO NECESITA
% ------------------------------------------------------------
% changepage: LOS BLOQUES A ANCHO COMPLETO (adjustwidth).
% DESDE EL CONTRATO 2 NO SE CARGA textpos: LA COLUMNA DERECHA YA NO
% VA EN POSICIÓN ABSOLUTA (4.4).
% ============================================================
\usepackage{changepage}

% ============================================================
% 3. ANCHO COMPLETO
% ------------------------------------------------------------
% EL TEXTO DEL ARTÍCULO DEJA LIBRE UNA COLUMNA AL MARGEN
% (marginparwidth). UN CUADRO, UNA FIGURA O UN TÍTULO A ANCHO
% COMPLETO LA OCUPAN: SE ENSANCHAN HACIA LA DERECHA LO QUE MIDEN EL
% MARGEN Y SU SEPARACIÓN. EL ANCHO LO DECIDE LA ESTÉTICA (geometry);
% ESTA CAPA SOLO GARANTIZA LOS NOMBRES QUE EMITE LA HOJA.
%   \fullwidthblock{…}       jats-to-latex.xsl
%   gbcuadrototal            tex-comun.xsl (CUADRO total: sí, SC-66)
% ============================================================
\newcommand{\fullwidthblock}[1]{%
  \begin{adjustwidth}{0pt}{-\dimexpr\marginparsep+\marginparwidth\relax}%
  \centering #1%
  \end{adjustwidth}%
}
\newenvironment{gbcuadrototal}
  {\begin{adjustwidth}{0pt}{-\dimexpr\marginparsep+\marginparwidth\relax}}
  {\end{adjustwidth}}

% ============================================================
% 4. PRIMERA PÁGINA (CONTRATO 2, A1.1b)
% ------------------------------------------------------------
% LO QUE EMITE jats-a-latex-articulo.xsl PARA LA PRIMERA PÁGINA Y
% PARA EL BLOQUE DE DECLARACIONES. ESTA CAPA DA LA MECÁNICA (CAJA,
% COLOCACIÓN, MEDIDA) Y UN ASPECTO NEUTRO; EL DE LA REVISTA LO PONE
% LA ESTÉTICA REDEFINIENDO LOS GANCHOS DE 4.2 CON \renewcommand (Y
% \tcbset{gbresumen/.style=…}).
% ============================================================

% ------------------------------------------------------------
% 4.0 IDIOMAS POR CÓDIGO
% ------------------------------------------------------------
% LA HOJA EMITE EL CÓDIGO DEL CANÓNICO (es, en, pt-BR) Y EL NOMBRE DE
% babel SALE DE \gbbabel@<código> (preambulo-numero.tex, DEL CATÁLOGO
% DE IDIOMAS DE LA BASE): ES EL MECANISMO DE \gbextranjero EN LIBROS.
% \selectlanguage Y otherlanguage NO EXPANDEN UN \csname EN SU
% ARGUMENTO (MEDIDO: TOMAN «csnamegbbabel@de…» COMO NOMBRE); POR ESO
% EL \edef, QUE LES PASA EL NOMBRE YA EXPANDIDO. UN CÓDIGO SIN NOMBRE
% FRENA CON UN MENSAJE Y NO CAMBIA DE IDIOMA: COMPONER CON OTRA
% PARTICIÓN ES UN ERROR QUE NADIE VE (MEDIDO, babel DE TeX Live 2023).
% ------------------------------------------------------------
\makeatletter
\newcommand*{\gbidiomadefinido}[1]{%
  \ifcsname gbbabel@#1\endcsname
    \expandafter\@firstoftwo
  \else
    \PackageError{gbpublisher}{El idioma «#1» no tiene nombre de babel}{%
      preambulo-numero.tex no lo define: hay que generar el PDF desde
      gbpublisher, que lo toma del catalogo de idiomas.}%
    \expandafter\@secondoftwo
  \fi}
\newcommand*{\gbselectlanguage}[1]{%
  \gbidiomadefinido{#1}{%
    \edef\gb@idioma{\noexpand\selectlanguage{\csname gbbabel@#1\endcsname}}%
    \gb@idioma}{}}
% EL ENTORNO DE OTRO IDIOMA: SI EL CÓDIGO NO TIENE NOMBRE, SIGUE EN EL
% IDIOMA DEL ARTÍCULO (YA FRENÓ CON EL MENSAJE)
\newenvironment{gbotroidioma}[1]
  {\gbidiomadefinido{#1}{%
     \edef\gb@idioma{\noexpand\otherlanguage{\csname gbbabel@#1\endcsname}}}%
     {\def\gb@idioma{\otherlanguage{\languagename}}}%
   \gb@idioma}
  {\endotherlanguage}
\makeatother

% ------------------------------------------------------------
% 4.1 LOS DATOS DEL ARTÍCULO
% ------------------------------------------------------------
% LA HOJA LOS REDEFINE AL EMPEZAR EL CUERPO Y LA ESTÉTICA LOS USA EN
% LA CABECERA Y EL PIE. ACÁ SE DECLARAN VACÍOS PARA QUE LA ESTÉTICA
% PUEDA NOMBRARLOS.
%   \articuloautorescorto  «Pérez», «Pérez y Gómez», «Pérez et al.»:
%                          EL MEMBRETE DE LAS INTERIORES (REDALYC 4.5)
%   \articuloeid           EL NÚMERO DE ARTÍCULO (MODELO continua)
% ------------------------------------------------------------
\newcommand{\articulotipo}{}
\newcommand{\articulotitulo}{}
\newcommand{\articulosubtitulo}{}
\newcommand{\articulotitulotrans}{}
\newcommand{\articuloautores}{}
\newcommand{\articuloautorescorto}{}
\newcommand{\articulodoi}{}
\newcommand{\articuloeid}{}
\newcommand{\articuloidiomaxml}{es}
\newcommand{\articulorevista}{}
\newcommand{\articuloissne}{}
\newcommand{\articulovolumen}{}
\newcommand{\articulonumero}{}
\newcommand{\articulofecha}{}
\newcommand{\articulopagina}{0}

% ------------------------------------------------------------
% 4.2 GANCHOS DE ASPECTO (LA ESTÉTICA LOS REDEFINE)
% ------------------------------------------------------------
\newcommand{\gbestilotitulo}{\sffamily\bfseries\Large}
\newcommand{\gbestilosubtitulo}{\sffamily\large}
\newcommand{\gbestilotitulotrans}{\sffamily\itshape\normalsize}
\newcommand{\gbestilocolumna}{\sffamily\small\RaggedRight}
\newcommand{\gbestilorotulo}{\bfseries}
\newcommand{\gbestiloautor}{\bfseries}
\newcommand{\gbestilodato}{\footnotesize}
\newcommand{\gbestiloresumen}{\sffamily\small}
\tcbset{gbresumen/.style={colback=white,colframe=black,boxrule=0.4pt,
  arc=0pt,left=4pt,right=4pt,top=3pt,bottom=3pt,
  before skip=4pt,after skip=4pt}}

% ------------------------------------------------------------
% 4.3 TÍTULOS A ANCHO COMPLETO: EL TEXTO MÁS LA COLUMNA DEL MARGEN
% ------------------------------------------------------------
\newlength{\gbanchototal}
\AtBeginDocument{\setlength{\gbanchototal}
  {\dimexpr\textwidth+\marginparsep+\marginparwidth\relax}}
\newcommand{\gbbloquetitulo}[3]{%
  \par\noindent\begin{minipage}[t]{\gbanchototal}#1#2\par\end{minipage}\par
  \vspace{#3}}
\newcommand{\gbtitulo}[1]{\gbbloquetitulo{\gbestilotitulo}{#1}{3mm}}
\newcommand{\gbsubtitulo}[1]{\gbbloquetitulo{\gbestilosubtitulo}{#1}{3mm}}
\newcommand{\gbtitulotrans}[1]{\gbbloquetitulo{\gbestilotitulotrans}{#1}{4mm}}

% ------------------------------------------------------------
% 4.4 LA COLUMNA DERECHA
% ------------------------------------------------------------
% gbcolumna GUARDA SU CONTENIDO EN UNA CAJA DEL ANCHO DEL MARGEN.
% \gbcolocarcolumna LA PONE EN LA COLUMNA DEL MARGEN CON SU BORDE
% DE ARRIBA EN EL PUNTO DONDE SE LLAMA (LA HOJA LO HACE ANTES DEL
% PRIMER RESUMEN), SIN OCUPAR LUGAR EN LA COLUMNA DEL TEXTO. ANTES
% LA MIDE CONTRA LO QUE QUEDA DE LA PÁGINA (\pagegoal-\pagetotal):
% SI NO ENTRA, COMPONE IGUAL Y DEJA EN EL .log UN AVISO
% «Package gbpublisher Warning: columna-primera-pagina», QUE EL
% SCRIPT TOMA PARA SALIR CON CÓDIGO 3 (DECISIÓN DE ALBERTO, 7/11).
% \fullcite VA EN UNA refsection PROPIA: EL ARTÍCULO NO SALE EN SUS
% PROPIAS REFERENCIAS. EL .bib LO ESCRIBE LA HOJA Y biber LO BUSCA
% POR BIBINPUTS.
% ------------------------------------------------------------
\newsavebox{\gbcajacolumna}
\newlength{\gbaltocolumna}
\newlength{\gbrestopagina}
\newenvironment{gbcolumna}
  {\begin{lrbox}{\gbcajacolumna}\begin{minipage}[t]{\marginparwidth}%
   \gbestilocolumna\setlength{\parindent}{0pt}\setlength{\parskip}{1pt}}
  {\end{minipage}\end{lrbox}%
   % lrbox GUARDA LA CAJA EN EL GRUPO DE ESTE ENTORNO: SIN \global SE
   % PIERDE EN EL \end{gbcolumna} (MEDIDO: LA COLUMNA NO SALÍA)
   \global\setbox\gbcajacolumna\box\gbcajacolumna}
\newcommand{\gbcolocarcolumna}{%
  \par
  \setlength{\gbaltocolumna}{\dimexpr\ht\gbcajacolumna+\dp\gbcajacolumna\relax}%
  \setlength{\gbrestopagina}{\dimexpr\pagegoal-\pagetotal\relax}%
  \ifdim\gbaltocolumna>\gbrestopagina
    \PackageWarningNoLine{gbpublisher}{columna-primera-pagina: la
      columna derecha mide \the\gbaltocolumna\space y en la pagina
      quedan \the\gbrestopagina}%
  \fi
  \nointerlineskip
  \vbox to 0pt{\hbox{\hskip\dimexpr\textwidth+\marginparsep\relax
    \usebox{\gbcajacolumna}}\vss}%
  \nointerlineskip}
\newcommand{\gbbloque}{\par\vspace{6pt}}
\newcommand{\gbentreautores}{\par\vspace{4pt}}
\newcommand{\gbrotulo}[1]{{\gbestilorotulo #1}\par}
\newcommand{\gbautor}[1]{{\gbestiloautor #1}\par}
\newcommand{\gbnombre}[1]{#1\par}
\newcommand{\gborcid}[1]{{\gbestilodato\url{#1}}\par}
\newcommand{\gbafiliacion}[1]{{\gbestilodato #1}\par}
% #1 EL RÓTULO DE LA CORRESPONDENCIA, VACÍO EN LOS DEMÁS CORREOS
\newcommand{\gbcorreo}[2]{{\gbestilodato
  \ifx\relax#1\relax\else{\gbestilorotulo #1:}\space\fi
  \href{mailto:#2}{\nolinkurl{#2}}}\par}
\newcommand{\gbdato}[2]{{\gbestilorotulo #1:}\space #2\par}
\newcommand{\gbtexto}[1]{{\gbestilodato #1}\par}
\newcommand{\gbcomocitar}[2]{{\gbestilodato
  \begin{refsection}[#1]\fullcite{#2}\end{refsection}}\par}
% COMPARTIR: \gbcompartir{\gbicono{ARCHIVO}{ENLACE}…}. CADA ÍCONO ES UN
% PNG (.hidden/fonts) DE ALTO \gbaltoicono, ENLAZADO; ENTRE UNO Y OTRO,
% UN ESPACIO DONDE LA LÍNEA PUEDE CORTAR
\newcommand{\gbaltoicono}{11pt}
\newcommand{\gbicono}[2]{\href{#2}{\includegraphics[height=\gbaltoicono]{#1}}\hspace{0.5em plus 0.2em}\allowbreak}
\newcommand{\gbcompartir}[1]{{\gbestilodato #1}\par}

% ------------------------------------------------------------
% 4.5 RESÚMENES
% ------------------------------------------------------------
% gbresumen{CÓDIGO DE IDIOMA}{RÓTULO}[estructurado]. EL RÓTULO VA EN
% LÍNEA; EN EL ESTRUCTURADO, SOLO EN SU LÍNEA, Y CADA SECCIÓN ABRE
% CON \gbresumensec{TÍTULO}, EN LÍNEA.
% ------------------------------------------------------------
\NewDocumentEnvironment{gbresumen}{m m o}
  {\begin{tcolorbox}[gbresumen]\begin{gbotroidioma}{#1}%
   \gbestiloresumen\setlength{\parindent}{0pt}%
   \IfNoValueTF{#3}{{\gbestilorotulo #2:}\space\ignorespaces}
                   {{\gbestilorotulo #2}\par\smallskip}}
  {\end{gbotroidioma}\end{tcolorbox}}
\newcommand{\gbresumensec}[1]{{\gbestilorotulo #1.}\space\ignorespaces}
\newcommand{\gbpalabrasclave}[2]{\par\smallskip{\gbestilorotulo #1:}\space #2}

% ------------------------------------------------------------
% 4.6 DECLARACIONES (AL FINAL, ANTES DE LAS REFERENCIAS)
% ------------------------------------------------------------
\newcommand{\gbdeclaraciones}[1]{\section*{#1}}
\newcommand{\gbdeclaracion}[1]{\par\noindent{\gbestilorotulo #1.}\space\ignorespaces}

% ============================================================
% 5. EL PDF DEL NÚMERO (CONTRATO 4, B5)
% ------------------------------------------------------------
% latex/maestro-numero.tex (numero-maestro-latex.xsl) CARGA LAS
% MISMAS CAPAS QUE UN ARTÍCULO, COMPONE LAS PRELIMINARES QUE EL
% NÚMERO DECLARA Y DESPUÉS INCLUYE EL PDF DE CADA ARTÍCULO, YA
% COMPUESTO Y VALIDADO, CON pdfpages (DECISIONES DE ALBERTO, 07/11 Y
% 12/11/2026):
%   - pdfpages NO CONSERVA LOS ENLACES DE LOS PDF QUE INCLUYE (MEDIDO
%     EL 07/11); pax, QUE LOS RECUPERA, NO FUNCIONA CON LuaLaTeX
%     (MEDIDO EL 12/11: «pdfTeX is too old»). LOS ENLACES QUEDAN EN
%     LOS PDF SUELTOS DE CADA ARTÍCULO;
%   - EL SUMARIO SÍ ENLAZA: CADA ARTÍCULO SE INCLUYE CON link Y
%     linkname, QUE DEJAN UN DESTINO <linkname>.<página> POR PÁGINA;
%   - LAS PRELIMINARES NO LLEVAN FOLIO; EL PRIMER ARTÍCULO EMPIEZA EN
%     LA PÁGINA QUE SIGUE AL SUMARIO, SIN BLANCAS;
%   - EL ORDEN LO DA EL GENERADOR (fm-0N, SC-44): ESTA CAPA SOLO
%     COMPONE.
% ESTA CAPA DA LA MECÁNICA Y UN ASPECTO NEUTRO; EL DE LA REVISTA LO
% PONE LA ESTÉTICA REDEFINIENDO LOS GANCHOS DE 5.1 CON \renewcommand.
% pdfpages LO CARGA EL MAESTRO: UN ARTÍCULO NO LO NECESITA.
% ============================================================
\makeatletter

% ------------------------------------------------------------
% 5.1 GANCHOS DE ASPECTO Y RÓTULOS (LA ESTÉTICA LOS REDEFINE)
% ------------------------------------------------------------
\newcommand{\gbestilosumariotitulo}{\sffamily\bfseries\Large}
\newcommand{\gbestilosumarioseccion}{\sffamily\bfseries\MakeUppercase}
\newcommand{\gbestilosumarioarticulo}{\bfseries}
\newcommand{\gbestilosumarioautores}{\itshape}
\newcommand{\gbestilolegal}{\sffamily\small\RaggedRight}
\newcommand{\gbestilolegalrevista}{\sffamily\bfseries\large}
\newcommand{\gbestilolegalrotulo}{\bfseries}
% LOS RÓTULOS FIJOS, EN CASTELLANO: UNA REVISTA EN OTRO IDIOMA LOS
% REDEFINE EN SU ESTÉTICA
\newcommand{\gbrotulosumario}{Sumario}
\newcommand{\gbrotulotapa}{Tapa}
\newcommand{\gbrotulolegal}{Página legal}
\newcommand{\gbrotulovolumen}{Vol.}
\newcommand{\gbrotulonumero}{n.\textsuperscript{o}}
\newcommand{\gbrotulosuplemento}{supl.}
\newcommand{\gbnombremes}[1]{\ifcase#1\or enero\or febrero\or marzo\or
  abril\or mayo\or junio\or julio\or agosto\or septiembre\or octubre\or
  noviembre\or diciembre\fi}

% ------------------------------------------------------------
% 5.2 LOS DATOS DEL NÚMERO
% ------------------------------------------------------------
% \gbdatosnumero{REVISTA}{RÓTULO DEL VOLUMEN}{VOLUMEN}{NÚMERO}
%               {SUPLEMENTO}{MES}{AÑO}{TÍTULO ESPECIAL}
% LOS GUARDA PARA LA PÁGINA LEGAL Y EL SUMARIO. CUALQUIERA PUEDE IR
% VACÍO. EL MES ES EL NÚMERO (1 A 12). EL RÓTULO VACÍO ES
% \gbrotulovolumen (revista.rotulo_volumen: «Año», «Tomo»…).
% \gbidentificacionnumero COMPONE «Vol. 3, n.º 2, noviembre de 2026»
% CON LO QUE HAYA.
% ------------------------------------------------------------
\newcommand{\gbdatosnumero}[8]{%
  \def\gbnumerorevista{#1}\def\gbnumerorotulo{#2}%
  \def\gbnumerovolumen{#3}\def\gbnumeronumero{#4}%
  \def\gbnumerosuplemento{#5}\def\gbnumeromes{#6}%
  \def\gbnumeroanio{#7}\def\gbnumerotituloespecial{#8}}
\gbdatosnumero{}{}{}{}{}{}{}{}
\newcommand{\gb@sep}{}
% \nobreakspace Y NO ~: CON babel-spanish ~ ES UN ATAJO ACTIVO QUE TOMA
% EL SIGUIENTE TOKEN (MEDIDO: «~de~\fi» SE COMÍA EL \fi)
\newcommand{\gbidentificacionnumero}{%
  \def\gb@sep{}%
  \ifx\gbnumerovolumen\@empty\else
    \ifx\gbnumerorotulo\@empty\gbrotulovolumen\else\gbnumerorotulo\fi
    \nobreakspace\gbnumerovolumen\def\gb@sep{, }\fi
  \ifx\gbnumeronumero\@empty\else\gb@sep\gbrotulonumero\nobreakspace
    \gbnumeronumero\def\gb@sep{, }\fi
  \ifx\gbnumerosuplemento\@empty\else\gb@sep\gbrotulosuplemento
    \nobreakspace\gbnumerosuplemento\def\gb@sep{, }\fi
  \ifx\gbnumeroanio\@empty\else\gb@sep
    \ifx\gbnumeromes\@empty\else\gbnombremes{\gbnumeromes}\space de\space\fi
    \gbnumeroanio\fi}

% ------------------------------------------------------------
% 5.3 PÁGINAS DE LAS PRELIMINARES
% ------------------------------------------------------------
% CADA PRELIMINAR EMPIEZA PÁGINA, SIN CABECERA NI FOLIO, Y LLEVA SU
% MARCADOR, AL MISMO NIVEL QUE EL DE CADA ARTÍCULO (1, EL DE section:
% CON 0 LOS ARTÍCULOS QUEDABAN COLGADOS DEL SUMARIO). EL TEXTO OCUPA EL
% ANCHO COMPLETO: LA CAJA MÁS LA COLUMNA
% DEL MARGEN (COMO \fullwidthblock, 3).
%   \gbpreliminar{MARCADOR}{ANCLA}   ABRE LA PÁGINA
% ------------------------------------------------------------
\newcommand{\gbpreliminar}[2]{%
  \clearpage\thispagestyle{empty}%
  \pdfbookmark[1]{#1}{#2}}
\newenvironment{gbanchocompleto}
  {\begin{adjustwidth}{0pt}{-\dimexpr\marginparsep+\marginparwidth\relax}}
  {\end{adjustwidth}}

% ------------------------------------------------------------
% 5.4 TAPA
% ------------------------------------------------------------
% \gbtapanumero{RUTA DE LA IMAGEN}: LA IMAGEN DEL NÚMERO
% (numero.imagen_tapa, EN media/) A PÁGINA COMPLETA, ESCALADA SIN
% DEFORMAR Y CENTRADA (pdfpages ACEPTA IMÁGENES). SIN IMAGEN EL
% SCRIPT FRENA ANTES (DECISIÓN DE ALBERTO, 12/11).
% ------------------------------------------------------------
\newcommand{\gbtapanumero}[1]{%
  \clearpage
  \includepdf[pagecommand={\thispagestyle{empty}\pdfbookmark[1]{\gbrotulotapa}{gbtapa}}]{#1}}

% ------------------------------------------------------------
% 5.5 PÁGINA LEGAL
% ------------------------------------------------------------
% gbpaginalegal ENVUELVE LOS DATOS, QUE EL GENERADOR EMITE EN ESTE
% ORDEN Y SOLO SI HAY DATO:
%   \gblegalrevista                     REVISTA E IDENTIFICACIÓN
%   \gblegaleditorial{NOMBRE}{INSTITUCIÓN}
%   \gblegalcredito{CARGO}{NOMBRES}     UNO POR CARGO, LOS NOMBRES YA
%                                       UNIDOS (VIGENTES A LA FECHA)
%   \gblegalorgano{NOMBRE}{CANTIDAD}{MIEMBROS}
%   \gblegalmiembro{NOMBRE}{AFILIACIÓN} DENTRO DE MIEMBROS
%   \gblegalissnimpreso{ISSN}  \gblegalissnelectronico{ISSN}
%   \gblegaldireccion{TEXTO}            DIRECCIÓN, REGIÓN, CIUDAD, PAÍS
%   \gblegalportal{URL}  \gblegalcorreo{CORREO}
%   \gblegalleyenda{TEXTO}  \gblegalaviso{TEXTO}
% UN ÓRGANO DE UN SOLO MIEMBRO VA EN LÍNEA («Coordinadora editorial:
% Nombre, afiliación»); CON VARIOS, UNO POR RENGLÓN. LA CANTIDAD LLEGA
% COMO DATO PARA QUE LA ESTÉTICA PUEDA DECIDIR OTRA COSA.
% ------------------------------------------------------------
\newenvironment{gbpaginalegal}
  {\gbpreliminar{\gbrotulolegal}{gblegal}%
   \begin{gbanchocompleto}\gbestilolegal\setlength{\parindent}{0pt}%
   \setlength{\parskip}{4pt plus 1pt}}
  {\end{gbanchocompleto}}
\newcommand{\gblegalbloque}{\par\medskip}
\newcommand{\gblegalrevista}{%
  {\gbestilolegalrevista\gbnumerorevista}\par
  \gbidentificacionnumero\par
  \ifx\gbnumerotituloespecial\@empty\else{\itshape\gbnumerotituloespecial}\par\fi
  \gblegalbloque}
\newcommand{\gblegaleditorial}[2]{%
  {\gbestilolegalrotulo #1}\par
  \if\relax\detokenize{#2}\relax\else #2\par\fi
  \gblegalbloque}
\newcommand{\gblegalcredito}[2]{{\gbestilolegalrotulo #1:}\space #2\par}
\newcommand{\gblegalorgano}[3]{%
  \gblegalbloque
  \ifnum#2=1\relax
    {\gbestilolegalrotulo #1:}\space\let\gblegalmiembro\gblegalmiembroenlinea #3\par
  \else
    {\gbestilolegalrotulo #1}\par\let\gblegalmiembro\gblegalmiembroenrenglon #3%
  \fi}
\newcommand{\gblegalmiembroenlinea}[2]{#1\if\relax\detokenize{#2}\relax\else, #2\fi}
\newcommand{\gblegalmiembroenrenglon}[2]{#1\if\relax\detokenize{#2}\relax\else\space(#2)\fi\par}
\newcommand{\gblegalmiembro}[2]{\gblegalmiembroenrenglon{#1}{#2}}
\newcommand{\gblegalissnimpreso}[1]{ISSN:\space #1\par}
\newcommand{\gblegalissnelectronico}[1]{e-ISSN:\space #1\par}
\newcommand{\gblegaldireccion}[1]{\gblegalbloque #1\par}
\newcommand{\gblegalportal}[1]{\url{#1}\par}
\newcommand{\gblegalcorreo}[1]{\href{mailto:#1}{#1}\par}
\newcommand{\gblegalleyenda}[1]{\gblegalbloque #1\par}
\newcommand{\gblegalaviso}[1]{\gblegalbloque{\footnotesize #1}\par}

% ------------------------------------------------------------
% 5.6 SUMARIO
% ------------------------------------------------------------
% gbsumarionumero ENVUELVE LAS ENTRADAS:
%   \gbsumarioseccion{TÍTULO}  LA SECCIÓN O EL DOSSIER, CUANDO CAMBIA
%   \gbsumarioarticulo{ANCLA}{IDIOMA}{TÍTULO}{AUTORES}{FOLIO}
% ANCLA ES EL linkname DEL ARTÍCULO (5.7): EL TÍTULO ENLAZA A SU
% PRIMERA PÁGINA. IDIOMA ES EL CÓDIGO DEL ARTÍCULO (4.0), PARA LA
% PARTICIÓN DEL TÍTULO. FOLIO: LA PÁGINA INICIAL (por_numero,
% por_volumen), EL e-id (continua) O VACÍO (por_articulo).
% ------------------------------------------------------------
\newenvironment{gbsumarionumero}
  {\gbpreliminar{\gbrotulosumario}{gbsumario}%
   \begin{gbanchocompleto}\setlength{\parindent}{0pt}%
   {\gbestilosumariotitulo\gbrotulosumario}\par\smallskip
   \gbidentificacionnumero\par\bigskip}
  {\end{gbanchocompleto}}
\newcommand{\gbsumarioseccion}[1]{%
  \par\addvspace{12pt}{\gbestilosumarioseccion{#1}}\par\nopagebreak\smallskip}
\newcommand{\gbsumarioarticulo}[5]{%
  \par\addvspace{6pt}%
  \begin{gbotroidioma}{#2}%
  {\gbestilosumarioarticulo\hyperlink{#1.1}{#3}}\nobreak\hfill\nobreak
  \if\relax\detokenize{#5}\relax\else\mbox{\hspace{1em}#5}\fi\par
  \end{gbotroidioma}%
  \if\relax\detokenize{#4}\relax\else{\gbestilosumarioautores #4}\par\fi}

% ------------------------------------------------------------
% 5.7 LOS ARTÍCULOS
% ------------------------------------------------------------
% \gbarticulonumero{ANCLA}{MARCADOR}{PDF}: TODAS LAS PÁGINAS DEL PDF
% DEL ARTÍCULO, TAL CUAL (CADA UNA YA TRAE SU CABECERA Y SU FOLIO),
% CON UN DESTINO POR PÁGINA (<ANCLA>.<página>) Y UN MARCADOR CON EL
% TÍTULO EN LA PRIMERA.
% ------------------------------------------------------------
\newcommand{\gbarticulonumero}[3]{%
  \includepdf[pages=-,link,linkname=#1,
    pagecommand={\thispagestyle{empty}},
    addtotoc={1,section,1,{#2},#1-marcador}]{#3}}
\makeatother
